% ===================================================================
%  stil_handreichung.tex – gemeinsamer Stil für Handreichungen
%  Einbinden mit % ===================================================================
%  stil_handreichung.tex – gemeinsamer Stil für Handreichungen
%  Einbinden mit % ===================================================================
%  stil_handreichung.tex – gemeinsamer Stil für Handreichungen
%  Einbinden mit % ===================================================================
%  stil_handreichung.tex – gemeinsamer Stil für Handreichungen
%  Einbinden mit \input{stil_handreichung} direkt nach \documentclass.
%  Kompilieren mit pdflatex (für Lautschrift mit Unicode: xelatex,
%  dann die fontenc/lmodern-Zeilen durch fontspec ersetzen).
% ===================================================================
\usepackage[T1]{fontenc}
\usepackage{lmodern}
\usepackage[ngerman]{babel}
\usepackage[autostyle]{csquotes}
\usepackage[a4paper,margin=2.3cm]{geometry}
\usepackage{microtype}
\usepackage{amsmath,amssymb}
\usepackage{xcolor}
\usepackage{graphicx}
\usepackage{enumitem}
\usepackage{booktabs}
\usepackage{tabularx}
\usepackage{longtable}
\usepackage{multicol}
\usepackage{tikz}
\usepackage{listings}
\usepackage[most]{tcolorbox}
\usepackage{xurl}
\usepackage{hyperref}

% --- Farben (Schulfarben; hier zentral anpassen) --------------------
\definecolor{mkgblue}{HTML}{2563EB}
\definecolor{mkgred}{HTML}{DC2626}
\definecolor{mkggreen}{HTML}{16A34A}
\definecolor{mkgpurple}{HTML}{7C3AED}
\definecolor{mkgorange}{HTML}{EA580C}
\definecolor{codebg}{HTML}{F4F4F7}
\definecolor{codeframe}{HTML}{D0D0D8}

\hypersetup{colorlinks=true, linkcolor=mkgblue, urlcolor=mkgblue}

% --- Quelltext ------------------------------------------------------
% language=TeX statt [LaTeX]TeX: die LaTeX-Variante stört innerhalb
% von tcblisting die Schachbefehle von skak.
\lstdefinestyle{tex}{
  basicstyle=\ttfamily\small, columns=fullflexible, keepspaces=true,
  breaklines=true, language=TeX, texcsstyle=*\color{mkgblue},
  commentstyle=\color{black!50}\itshape,
  literate={ä}{{\"a}}1 {ö}{{\"o}}1 {ü}{{\"u}}1 {Ä}{{\"A}}1 {Ö}{{\"O}}1
           {Ü}{{\"U}}1 {ß}{{\ss}}1 {–}{{--}}1
}
\lstset{style=tex, backgroundcolor=\color{codebg}, frame=single,
        rulecolor=\color{codeframe}, framesep=4pt}

% --- Beispiel-Kästen: Code wird angezeigt UND ausgeführt ------------
% links Code, rechts Ergebnis (optional: Breite der rechten Spalte,
% z. B. \begin{beispiel}[7cm]). Das „!“ vor O verhindert, dass beim Suchen
% nach dem optionalen Argument der Zeilenumbruch übersprungen und der
% erste Befehl der nächsten Zeile verschluckt wird.
\DeclareTCBListing{beispiel}{ !O{5.6cm} }{enhanced, sharp corners, breakable=false,
  colback=codebg, colframe=codeframe, boxrule=0.5pt,
  listing side text, righthand width=#1, sidebyside gap=5mm,
  segmentation style={draw=codeframe, solid},
  bicolor, colbacklower=white, listing options={style=tex, frame=none, backgroundcolor={}}}
% oben Code, unten Ergebnis
\newtcblisting{beispielunten}{enhanced, sharp corners,
  colback=codebg, colframe=codeframe, boxrule=0.5pt,
  listing above text, colbacklower=white, bicolor,
  segmentation style={draw=codeframe, solid},
  listing options={style=tex, frame=none, backgroundcolor={}}}

% --- Merke-, Hinweis- und Tipp-Kästen -------------------------------
\newtcolorbox{merke}[1][Merke]{enhanced, breakable, sharp corners,
  colback=mkgorange!6, colframe=mkgorange, boxrule=0pt, leftrule=3pt,
  fonttitle=\bfseries\color{mkgorange}, coltitle=mkgorange,
  colbacktitle=mkgorange!6, title={#1}, attach title to upper={\ }}
\newtcolorbox{hinweis}[1][Achtung]{enhanced, breakable, sharp corners,
  colback=mkgred!5, colframe=mkgred, boxrule=0pt, leftrule=3pt,
  fonttitle=\bfseries, coltitle=mkgred, colbacktitle=mkgred!5,
  title={#1}, attach title to upper={\ }}
\newtcolorbox{tipp}[1][Tipp für den Unterricht]{enhanced, breakable, sharp corners,
  colback=mkggreen!6, colframe=mkggreen, boxrule=0pt, leftrule=3pt,
  fonttitle=\bfseries, coltitle=mkggreen, colbacktitle=mkggreen!6,
  title={#1}, attach title to upper={\ }}

% --- Kurzbefehle ----------------------------------------------------
\newcommand\befehl[1]{\texttt{\textcolor{mkgblue}{\textbackslash #1}}}
\newcommand\opt[1]{\texttt{\textcolor{mkgpurple}{#1}}}
\newcommand\paket[1]{\textsf{\bfseries #1}}

% --- Titelblock -----------------------------------------------------
% \Titel{Haupttitel}{Untertitel}
\newcommand\Titel[2]{%
  \begin{tcolorbox}[enhanced, sharp corners, colback=mkgblue, colframe=mkgblue,
    coltext=white, boxsep=6pt, left=8pt]
    {\small\sffamily HANDREICHUNG}\par\smallskip
    {\LARGE\bfseries\sffamily #1}\par\smallskip
    {\large\sffamily #2}
  \end{tcolorbox}\bigskip}

\setlength{\parindent}{0pt}
\setlength{\parskip}{0.5em}
\setlist{itemsep=2pt, topsep=4pt}
 direkt nach \documentclass.
%  Kompilieren mit pdflatex (für Lautschrift mit Unicode: xelatex,
%  dann die fontenc/lmodern-Zeilen durch fontspec ersetzen).
% ===================================================================
\usepackage[T1]{fontenc}
\usepackage{lmodern}
\usepackage[ngerman]{babel}
\usepackage[autostyle]{csquotes}
\usepackage[a4paper,margin=2.3cm]{geometry}
\usepackage{microtype}
\usepackage{amsmath,amssymb}
\usepackage{xcolor}
\usepackage{graphicx}
\usepackage{enumitem}
\usepackage{booktabs}
\usepackage{tabularx}
\usepackage{longtable}
\usepackage{multicol}
\usepackage{tikz}
\usepackage{listings}
\usepackage[most]{tcolorbox}
\usepackage{xurl}
\usepackage{hyperref}

% --- Farben (Schulfarben; hier zentral anpassen) --------------------
\definecolor{mkgblue}{HTML}{2563EB}
\definecolor{mkgred}{HTML}{DC2626}
\definecolor{mkggreen}{HTML}{16A34A}
\definecolor{mkgpurple}{HTML}{7C3AED}
\definecolor{mkgorange}{HTML}{EA580C}
\definecolor{codebg}{HTML}{F4F4F7}
\definecolor{codeframe}{HTML}{D0D0D8}

\hypersetup{colorlinks=true, linkcolor=mkgblue, urlcolor=mkgblue}

% --- Quelltext ------------------------------------------------------
% language=TeX statt [LaTeX]TeX: die LaTeX-Variante stört innerhalb
% von tcblisting die Schachbefehle von skak.
\lstdefinestyle{tex}{
  basicstyle=\ttfamily\small, columns=fullflexible, keepspaces=true,
  breaklines=true, language=TeX, texcsstyle=*\color{mkgblue},
  commentstyle=\color{black!50}\itshape,
  literate={ä}{{\"a}}1 {ö}{{\"o}}1 {ü}{{\"u}}1 {Ä}{{\"A}}1 {Ö}{{\"O}}1
           {Ü}{{\"U}}1 {ß}{{\ss}}1 {–}{{--}}1
}
\lstset{style=tex, backgroundcolor=\color{codebg}, frame=single,
        rulecolor=\color{codeframe}, framesep=4pt}

% --- Beispiel-Kästen: Code wird angezeigt UND ausgeführt ------------
% links Code, rechts Ergebnis (optional: Breite der rechten Spalte,
% z. B. \begin{beispiel}[7cm]). Das „!“ vor O verhindert, dass beim Suchen
% nach dem optionalen Argument der Zeilenumbruch übersprungen und der
% erste Befehl der nächsten Zeile verschluckt wird.
\DeclareTCBListing{beispiel}{ !O{5.6cm} }{enhanced, sharp corners, breakable=false,
  colback=codebg, colframe=codeframe, boxrule=0.5pt,
  listing side text, righthand width=#1, sidebyside gap=5mm,
  segmentation style={draw=codeframe, solid},
  bicolor, colbacklower=white, listing options={style=tex, frame=none, backgroundcolor={}}}
% oben Code, unten Ergebnis
\newtcblisting{beispielunten}{enhanced, sharp corners,
  colback=codebg, colframe=codeframe, boxrule=0.5pt,
  listing above text, colbacklower=white, bicolor,
  segmentation style={draw=codeframe, solid},
  listing options={style=tex, frame=none, backgroundcolor={}}}

% --- Merke-, Hinweis- und Tipp-Kästen -------------------------------
\newtcolorbox{merke}[1][Merke]{enhanced, breakable, sharp corners,
  colback=mkgorange!6, colframe=mkgorange, boxrule=0pt, leftrule=3pt,
  fonttitle=\bfseries\color{mkgorange}, coltitle=mkgorange,
  colbacktitle=mkgorange!6, title={#1}, attach title to upper={\ }}
\newtcolorbox{hinweis}[1][Achtung]{enhanced, breakable, sharp corners,
  colback=mkgred!5, colframe=mkgred, boxrule=0pt, leftrule=3pt,
  fonttitle=\bfseries, coltitle=mkgred, colbacktitle=mkgred!5,
  title={#1}, attach title to upper={\ }}
\newtcolorbox{tipp}[1][Tipp für den Unterricht]{enhanced, breakable, sharp corners,
  colback=mkggreen!6, colframe=mkggreen, boxrule=0pt, leftrule=3pt,
  fonttitle=\bfseries, coltitle=mkggreen, colbacktitle=mkggreen!6,
  title={#1}, attach title to upper={\ }}

% --- Kurzbefehle ----------------------------------------------------
\newcommand\befehl[1]{\texttt{\textcolor{mkgblue}{\textbackslash #1}}}
\newcommand\opt[1]{\texttt{\textcolor{mkgpurple}{#1}}}
\newcommand\paket[1]{\textsf{\bfseries #1}}

% --- Titelblock -----------------------------------------------------
% \Titel{Haupttitel}{Untertitel}
\newcommand\Titel[2]{%
  \begin{tcolorbox}[enhanced, sharp corners, colback=mkgblue, colframe=mkgblue,
    coltext=white, boxsep=6pt, left=8pt]
    {\small\sffamily HANDREICHUNG}\par\smallskip
    {\LARGE\bfseries\sffamily #1}\par\smallskip
    {\large\sffamily #2}
  \end{tcolorbox}\bigskip}

\setlength{\parindent}{0pt}
\setlength{\parskip}{0.5em}
\setlist{itemsep=2pt, topsep=4pt}
 direkt nach \documentclass.
%  Kompilieren mit pdflatex (für Lautschrift mit Unicode: xelatex,
%  dann die fontenc/lmodern-Zeilen durch fontspec ersetzen).
% ===================================================================
\usepackage[T1]{fontenc}
\usepackage{lmodern}
\usepackage[ngerman]{babel}
\usepackage[autostyle]{csquotes}
\usepackage[a4paper,margin=2.3cm]{geometry}
\usepackage{microtype}
\usepackage{amsmath,amssymb}
\usepackage{xcolor}
\usepackage{graphicx}
\usepackage{enumitem}
\usepackage{booktabs}
\usepackage{tabularx}
\usepackage{longtable}
\usepackage{multicol}
\usepackage{tikz}
\usepackage{listings}
\usepackage[most]{tcolorbox}
\usepackage{xurl}
\usepackage{hyperref}

% --- Farben (Schulfarben; hier zentral anpassen) --------------------
\definecolor{mkgblue}{HTML}{2563EB}
\definecolor{mkgred}{HTML}{DC2626}
\definecolor{mkggreen}{HTML}{16A34A}
\definecolor{mkgpurple}{HTML}{7C3AED}
\definecolor{mkgorange}{HTML}{EA580C}
\definecolor{codebg}{HTML}{F4F4F7}
\definecolor{codeframe}{HTML}{D0D0D8}

\hypersetup{colorlinks=true, linkcolor=mkgblue, urlcolor=mkgblue}

% --- Quelltext ------------------------------------------------------
% language=TeX statt [LaTeX]TeX: die LaTeX-Variante stört innerhalb
% von tcblisting die Schachbefehle von skak.
\lstdefinestyle{tex}{
  basicstyle=\ttfamily\small, columns=fullflexible, keepspaces=true,
  breaklines=true, language=TeX, texcsstyle=*\color{mkgblue},
  commentstyle=\color{black!50}\itshape,
  literate={ä}{{\"a}}1 {ö}{{\"o}}1 {ü}{{\"u}}1 {Ä}{{\"A}}1 {Ö}{{\"O}}1
           {Ü}{{\"U}}1 {ß}{{\ss}}1 {–}{{--}}1
}
\lstset{style=tex, backgroundcolor=\color{codebg}, frame=single,
        rulecolor=\color{codeframe}, framesep=4pt}

% --- Beispiel-Kästen: Code wird angezeigt UND ausgeführt ------------
% links Code, rechts Ergebnis (optional: Breite der rechten Spalte,
% z. B. \begin{beispiel}[7cm]). Das „!“ vor O verhindert, dass beim Suchen
% nach dem optionalen Argument der Zeilenumbruch übersprungen und der
% erste Befehl der nächsten Zeile verschluckt wird.
\DeclareTCBListing{beispiel}{ !O{5.6cm} }{enhanced, sharp corners, breakable=false,
  colback=codebg, colframe=codeframe, boxrule=0.5pt,
  listing side text, righthand width=#1, sidebyside gap=5mm,
  segmentation style={draw=codeframe, solid},
  bicolor, colbacklower=white, listing options={style=tex, frame=none, backgroundcolor={}}}
% oben Code, unten Ergebnis
\newtcblisting{beispielunten}{enhanced, sharp corners,
  colback=codebg, colframe=codeframe, boxrule=0.5pt,
  listing above text, colbacklower=white, bicolor,
  segmentation style={draw=codeframe, solid},
  listing options={style=tex, frame=none, backgroundcolor={}}}

% --- Merke-, Hinweis- und Tipp-Kästen -------------------------------
\newtcolorbox{merke}[1][Merke]{enhanced, breakable, sharp corners,
  colback=mkgorange!6, colframe=mkgorange, boxrule=0pt, leftrule=3pt,
  fonttitle=\bfseries\color{mkgorange}, coltitle=mkgorange,
  colbacktitle=mkgorange!6, title={#1}, attach title to upper={\ }}
\newtcolorbox{hinweis}[1][Achtung]{enhanced, breakable, sharp corners,
  colback=mkgred!5, colframe=mkgred, boxrule=0pt, leftrule=3pt,
  fonttitle=\bfseries, coltitle=mkgred, colbacktitle=mkgred!5,
  title={#1}, attach title to upper={\ }}
\newtcolorbox{tipp}[1][Tipp für den Unterricht]{enhanced, breakable, sharp corners,
  colback=mkggreen!6, colframe=mkggreen, boxrule=0pt, leftrule=3pt,
  fonttitle=\bfseries, coltitle=mkggreen, colbacktitle=mkggreen!6,
  title={#1}, attach title to upper={\ }}

% --- Kurzbefehle ----------------------------------------------------
\newcommand\befehl[1]{\texttt{\textcolor{mkgblue}{\textbackslash #1}}}
\newcommand\opt[1]{\texttt{\textcolor{mkgpurple}{#1}}}
\newcommand\paket[1]{\textsf{\bfseries #1}}

% --- Titelblock -----------------------------------------------------
% \Titel{Haupttitel}{Untertitel}
\newcommand\Titel[2]{%
  \begin{tcolorbox}[enhanced, sharp corners, colback=mkgblue, colframe=mkgblue,
    coltext=white, boxsep=6pt, left=8pt]
    {\small\sffamily HANDREICHUNG}\par\smallskip
    {\LARGE\bfseries\sffamily #1}\par\smallskip
    {\large\sffamily #2}
  \end{tcolorbox}\bigskip}

\setlength{\parindent}{0pt}
\setlength{\parskip}{0.5em}
\setlist{itemsep=2pt, topsep=4pt}
 direkt nach \documentclass.
%  Kompilieren mit pdflatex (für Lautschrift mit Unicode: xelatex,
%  dann die fontenc/lmodern-Zeilen durch fontspec ersetzen).
% ===================================================================
\usepackage[T1]{fontenc}
\usepackage{lmodern}
\usepackage[ngerman]{babel}
\usepackage[autostyle]{csquotes}
\usepackage[a4paper,margin=2.3cm]{geometry}
\usepackage{microtype}
\usepackage{amsmath,amssymb}
\usepackage{xcolor}
\usepackage{graphicx}
\usepackage{enumitem}
\usepackage{booktabs}
\usepackage{tabularx}
\usepackage{longtable}
\usepackage{multicol}
\usepackage{tikz}
\usepackage{listings}
\usepackage[most]{tcolorbox}
\usepackage{xurl}
\usepackage{hyperref}

% --- Farben (Schulfarben; hier zentral anpassen) --------------------
\definecolor{mkgblue}{HTML}{2563EB}
\definecolor{mkgred}{HTML}{DC2626}
\definecolor{mkggreen}{HTML}{16A34A}
\definecolor{mkgpurple}{HTML}{7C3AED}
\definecolor{mkgorange}{HTML}{EA580C}
\definecolor{codebg}{HTML}{F4F4F7}
\definecolor{codeframe}{HTML}{D0D0D8}

\hypersetup{colorlinks=true, linkcolor=mkgblue, urlcolor=mkgblue}

% --- Quelltext ------------------------------------------------------
% language=TeX statt [LaTeX]TeX: die LaTeX-Variante stört innerhalb
% von tcblisting die Schachbefehle von skak.
\lstdefinestyle{tex}{
  basicstyle=\ttfamily\small, columns=fullflexible, keepspaces=true,
  breaklines=true, language=TeX, texcsstyle=*\color{mkgblue},
  commentstyle=\color{black!50}\itshape,
  literate={ä}{{\"a}}1 {ö}{{\"o}}1 {ü}{{\"u}}1 {Ä}{{\"A}}1 {Ö}{{\"O}}1
           {Ü}{{\"U}}1 {ß}{{\ss}}1 {–}{{--}}1
}
\lstset{style=tex, backgroundcolor=\color{codebg}, frame=single,
        rulecolor=\color{codeframe}, framesep=4pt}

% --- Beispiel-Kästen: Code wird angezeigt UND ausgeführt ------------
% links Code, rechts Ergebnis (optional: Breite der rechten Spalte,
% z. B. \begin{beispiel}[7cm]). Das „!“ vor O verhindert, dass beim Suchen
% nach dem optionalen Argument der Zeilenumbruch übersprungen und der
% erste Befehl der nächsten Zeile verschluckt wird.
\DeclareTCBListing{beispiel}{ !O{5.6cm} }{enhanced, sharp corners, breakable=false,
  colback=codebg, colframe=codeframe, boxrule=0.5pt,
  listing side text, righthand width=#1, sidebyside gap=5mm,
  segmentation style={draw=codeframe, solid},
  bicolor, colbacklower=white, listing options={style=tex, frame=none, backgroundcolor={}}}
% oben Code, unten Ergebnis
\newtcblisting{beispielunten}{enhanced, sharp corners,
  colback=codebg, colframe=codeframe, boxrule=0.5pt,
  listing above text, colbacklower=white, bicolor,
  segmentation style={draw=codeframe, solid},
  listing options={style=tex, frame=none, backgroundcolor={}}}

% --- Merke-, Hinweis- und Tipp-Kästen -------------------------------
\newtcolorbox{merke}[1][Merke]{enhanced, breakable, sharp corners,
  colback=mkgorange!6, colframe=mkgorange, boxrule=0pt, leftrule=3pt,
  fonttitle=\bfseries\color{mkgorange}, coltitle=mkgorange,
  colbacktitle=mkgorange!6, title={#1}, attach title to upper={\ }}
\newtcolorbox{hinweis}[1][Achtung]{enhanced, breakable, sharp corners,
  colback=mkgred!5, colframe=mkgred, boxrule=0pt, leftrule=3pt,
  fonttitle=\bfseries, coltitle=mkgred, colbacktitle=mkgred!5,
  title={#1}, attach title to upper={\ }}
\newtcolorbox{tipp}[1][Tipp für den Unterricht]{enhanced, breakable, sharp corners,
  colback=mkggreen!6, colframe=mkggreen, boxrule=0pt, leftrule=3pt,
  fonttitle=\bfseries, coltitle=mkggreen, colbacktitle=mkggreen!6,
  title={#1}, attach title to upper={\ }}

% --- Kurzbefehle ----------------------------------------------------
\newcommand\befehl[1]{\texttt{\textcolor{mkgblue}{\textbackslash #1}}}
\newcommand\opt[1]{\texttt{\textcolor{mkgpurple}{#1}}}
\newcommand\paket[1]{\textsf{\bfseries #1}}

% --- Titelblock -----------------------------------------------------
% \Titel{Haupttitel}{Untertitel}
\newcommand\Titel[2]{%
  \begin{tcolorbox}[enhanced, sharp corners, colback=mkgblue, colframe=mkgblue,
    coltext=white, boxsep=6pt, left=8pt]
    {\small\sffamily HANDREICHUNG}\par\smallskip
    {\LARGE\bfseries\sffamily #1}\par\smallskip
    {\large\sffamily #2}
  \end{tcolorbox}\bigskip}

\setlength{\parindent}{0pt}
\setlength{\parskip}{0.5em}
\setlist{itemsep=2pt, topsep=4pt}
